% \iffalse
%<*driver>
\documentclass{ltxdoc}
\usepackage[T2A]{fontenc}
\usepackage[utf8]{inputenc}
\usepackage[russian]{babel}
\EnableCrossrefs
\CodelineIndex
\RecordChanges
\begin{document}
  \DocInput{poly-report.dtx}
\end{document}
%</driver>
% \fi
%
% \GetFileInfo{poly-report.sty}
% \title{Пакет \textsf{poly-report}}
% \author{Ilya Semenov}
% \date{\filedate}
% \maketitle
%
% Общие настройки и макросы для учебных отчётов в \LaTeX.
%
% \StopEventually{}
%
% \section{Реализация}
%
%    \begin{macrocode}
%<*package>
\NeedsTeXFormat{LaTeX2e}
\ProvidesPackage{poly-report}[2026/09/25 Shared university report style]

\RequirePackage[14pt]{extsizes}
\RequirePackage[T2A]{fontenc}
\RequirePackage[utf8]{inputenc}
\RequirePackage[russian]{babel}
\RequirePackage{geometry}
\geometry{left=2.5cm, right=2.5cm, top=2cm, bottom=2cm}
\RequirePackage[small]{titlesec}
\RequirePackage{enumitem}

\RequirePackage{xcolor}
\RequirePackage{tikz}
\usetikzlibrary{shapes.geometric}
\usetikzlibrary{calc}
\usetikzlibrary{positioning}
\usetikzlibrary{patterns.meta}
\usetikzlibrary{arrows}
\usetikzlibrary{fit}

\RequirePackage{amsmath}
\RequirePackage{amssymb}
\RequirePackage{caption}
\RequirePackage{multicol}
\RequirePackage{float}
\floatstyle{plaintop}
\restylefloat{listing}

\RequirePackage{tabularx}
\RequirePackage{tabularray}
\DeclareTblrTemplate{contfoot-text}{normal}{Продолжение на следующей странице}
\SetTblrTemplate{contfoot-text}{normal}
\DeclareTblrTemplate{conthead-text}{normal}{(Продолжение)}
\SetTblrTemplate{conthead-text}{normal}
\UseTblrLibrary{amsmath}
\UseTblrLibrary{tikz}

\RequirePackage{booktabs}
\RequirePackage{multirow}
\RequirePackage{indentfirst}
\RequirePackage{pdfpages}
\setlength{\columnsep}{0.8cm}

\captionsetup[listing]{
    position=above,
    singlelinecheck=false,
    justification=raggedright,
    skip=0pt,
    font=normalsize,
}

\RequirePackage{fvextra}
\RequirePackage{minted}
\usemintedstyle{lovelace}
\setminted{
    fontsize=\small,
    baselinestretch=1.0,
    breaklines,
    frame=none,
    numbersep=8pt,
}
\renewcommand{\theFancyVerbLine}{%
    \small\ttfamily\textcolor{gray}{\arabic{FancyVerbLine}}%
}
\renewcommand{\listingscaption}{Листинг}

\RequirePackage[
    colorlinks=true,
    linkcolor=blue,
    urlcolor=blue,
    citecolor=blue,
    filecolor=blue,
    pdfborder={0 0 0}
]{hyperref}

% \img[width][includegraphics options]{path}{caption}{post-caption commands}
\NewDocumentCommand{\img}{O{0.9} O{} m m m}{
    \begin{figure}[H]
        \centering
        \includegraphics[width=#1\textwidth,#2]{#3}
        \caption{#4}
        #5
    \end{figure}
}

\let\screenshot\img

% \imgGrid[gap](colspec){cells}{caption}{post-caption commands}
\NewDocumentCommand{\imgGrid}{O{1ex} D(){X[c,m] X[c,m]} m m m}{
    \begin{figure}[H]
        \centering
        \begin{tblr}{
            width=\textwidth,
            colspec={#2},
            columns={colsep=\dimexpr#1/2\relax},
            column{1}={leftsep=0pt},
            column{Z}={rightsep=0pt},
            rows={rowsep=\dimexpr#1/2\relax},
            row{1}={abovesep=0pt},
            row{Z}={belowsep=0pt}
        }
            #3
        \end{tblr}
        \caption{#4}
        #5
    \end{figure}
}

\NewDocumentCommand{\gridImage}{O{} m}{
    \includegraphics[width=\linewidth,#1]{#2}
}

% Semantic anchor on an included PDF page.
% \exerciseAnchor{name}{x from page left}{y from page top}
\NewDocumentCommand{\exerciseAnchor}{m m m}{
    \coordinate (#1) at
        ([xshift=#2,yshift=-#3]current page.north west);
}

% Place a screenshot relative to a semantic PDF-page anchor.
% \screenAt[includegraphics options]{anchor}{x offset}{y offset}{width}{path}
\NewDocumentCommand{\screenAt}{O{} m m m m m}{
    \node[
        anchor=north west,
        inner sep=0pt,
        xshift=#3,
        yshift=-#4
    ] at (#2) {
        \includegraphics[width=#5,#1]{#6}
    };
}

% \codeListing[language][minted options]{path}{caption}{post-caption commands}
\NewDocumentCommand{\codeListing}{O{cpp} O{} m m m}{
    \par\addvspace{\baselineskip}
    \captionof{listing}{#4}
    #5
    \inputminted[
        linenos,
        firstnumber=1,
        numbersep=10pt,
        fontsize=\small,
        #2
    ]{#1}{#3}
    \par\addvspace{\baselineskip}
}

% \fullPageFigure{path}{caption}{post-caption commands}
\newcommand{\fullPageFigure}[3]{
    \clearpage
    \begingroup
        \pdfpagewidth=420mm
        \pdfpageheight=297mm
        \paperwidth=\pdfpagewidth
        \paperheight=\pdfpageheight
        \textwidth=\dimexpr\paperwidth-30mm\relax
        \textheight=\dimexpr\paperheight-30mm\relax
        \hsize=\textwidth
        \linewidth=\textwidth
        \columnwidth=\textwidth
        \oddsidemargin=\dimexpr15mm-1in\relax
        \evensidemargin=\oddsidemargin
        \topmargin=\dimexpr15mm-1in\relax
        \headheight=0pt
        \headsep=0pt
        \footskip=0pt
        \thispagestyle{empty}
        \null
        \refstepcounter{figure}
        #3
        \begin{tikzpicture}[remember picture, overlay]
            \node[anchor=north] at ([yshift=-15mm]current page.north) {
                \includegraphics[
                    width=390mm,
                    height=235mm,
                    keepaspectratio
                ]{#1}
            };
            \node[
                anchor=south,
                align=center,
                font=\normalsize
            ] at ([yshift=10mm]current page.south) {
                Рис.~\thefigure: #2\\[0.45cm]
                \thepage
            };
        \end{tikzpicture}
        \clearpage
    \endgroup
    \restoregeometry
    \pdfpagewidth=\paperwidth
    \pdfpageheight=\paperheight
}

\newcolumntype{Y}{>{\centering\arraybackslash}X}

% Entity-relationship diagram styles derived from tikz-er2 by Pavel Caldo.
\tikzstyle{every entity} = []
\tikzstyle{every weak entity} = []
\tikzstyle{every attribute} = []
\tikzstyle{every relationship} = []
\tikzstyle{every link} = []
\tikzstyle{every isa} = []
\tikzstyle{link} = [>=triangle 60, draw, thick, every link]
\tikzstyle{total} = [link, double, double distance=3pt]
\tikzstyle{entity} = [rectangle, draw, black, very thick,
                      minimum width=6em, minimum height=3em,
                      every entity]
\tikzstyle{weak entity} = [entity, double, double distance=2pt,
                           every weak entity]
\tikzstyle{ident relationship} = [relationship, double, double distance=2pt]
\tikzstyle{attribute} = [ellipse, draw, thin, black,
                         minimum width=4em, minimum height=2em,
                         text width=4em, align=center,
                         every attribute]
\tikzstyle{multi attribute} = [attribute, double, double distance=1.5pt]
\tikzstyle{derived attribute} = [attribute, dashed]
\tikzstyle{key attribute} = [attribute]
\tikzstyle{relationship} = [diamond, draw, black, very thick,
                            minimum width=6em, minimum height=4em,
                            aspect=1.5,
                            every relationship]
\tikzstyle{isa} = [regular polygon, regular polygon sides=3,
                   draw, black, very thick,
                   minimum width=4em, minimum height=4em,
                   every isa]
\providecommand{\key}[1]{\underline{#1}}
%</package>
%    \end{macrocode}
%
% \Finale
